\documentclass[11pt]{article}
\ifdefined\pdfinfoomitdate\pdfinfoomitdate=1\fi
\ifdefined\pdftrailerid\pdftrailerid{}\fi
\usepackage[a4paper,hmargin=26mm,vmargin=25mm]{geometry}
\usepackage[T1]{fontenc}
\usepackage{lmodern,microtype,amsmath,amssymb,amsthm,mathtools}
\usepackage[hidelinks,bookmarksnumbered]{hyperref}
\hypersetup{pdftitle={Natural boundary and differential transcendence of the isotropic square-lattice Ising susceptibility},pdfauthor={Dehao Lin}}
\newtheorem*{theorem}{Theorem}
\newtheorem{lemma}{Lemma}
\newtheorem{proposition}[lemma]{Proposition}
\newcommand{\ii}{\mathrm i}
\newcommand{\dd}{\,\mathrm d}
\newcommand{\C}{\mathbb C}
\newcommand{\R}{\mathbb R}
\newcommand{\Z}{\mathbb Z}
\newcommand{\T}{\mathbb T}
\DeclareMathOperator{\arsinh}{arsinh}
\setlength{\emergencystretch}{1.5em}
\raggedbottom
\title{Natural boundary and differential transcendence of the\\
isotropic square-lattice Ising susceptibility}
\author{Dehao Lin}\date{}
\makeatletter
\renewcommand{\@maketitle}{\begin{center}
  \vspace*{.5\baselineskip}{\LARGE\@title\par}\vspace{.5\baselineskip}
  {\large\@author\par}
  {\small School of Physics, Sun Yat-sen University\par
  \href{https://orcid.org/0009-0001-4551-8490}{ORCID 0009-0001-4551-8490}\par}
\end{center}}
\makeatother
\begin{document}
\maketitle
\begin{abstract}
For the zero-field isotropic Ising model on the square lattice, we prove
that the unit circle in $s=\sinh(2\beta J)$ is a natural boundary for the
full susceptibility in each phase, and that both susceptibilities are
differentially transcendental over $\C(s)$. At a dense set of boundary
points we identify the first particle order with a singular contribution.
All collisions at that order have the same phase. A Gaussian estimate
on the lattice controls the derivatives of the infinite tail, so the
singularity persists in the full sum. The order of the first divergent
derivative becomes arbitrarily large on every boundary arc. Together
with the continuous boundary values of the lower derivatives, this
rules out an algebraic differential equation.
\end{abstract}

\section{The susceptibility in the two phases}

Guttmann and Enting found evidence of accumulating singularities in the
anisotropic high-temperature series \cite{GE}.
Nickel found singularities of the odd-particle integrals of the
isotropic square-lattice Ising susceptibility at a dense set of points
on the unit circle. He proposed that the circle is a natural boundary
of the full susceptibility \cite{Nickel}. The passage from a single
particle integral to the infinite sum requires control of cancellation.
Guttmann, Jensen, Maillard and Pantone asked whether the full
susceptibility satisfies an algebraic differential equation and tested
the series on both sides of the critical temperature \cite[Sec.~6]{GJMP}.

Let $J>0$ be the nearest-neighbor coupling, $\beta$ the inverse
temperature and $s=\sinh(2\beta J)$. With spins of modulus one, the
reduced zero-field susceptibilities are
\[
 X_H(s)=\sum_{a,b\in\Z}\langle\sigma_{00}\sigma_{ab}\rangle,\qquad
 X_L(s)=\sum_{a,b\in\Z}
 \bigl(\langle\sigma_{00}\sigma_{ab}\rangle_+-M^2\bigr),
 \qquad M^2=(1-s^{-4})^{1/4}.
\]
The first sum is taken at $0<s<1$, the second at $s>1$ in the plus
pure state. Spin reversal gives the same connected sum in the minus
state. Both sums include the origin. For a magnetic field in energy
units the physical susceptibilities are $\chi_\nu=\beta X_\nu$,
with the ordered-state response taken from positive field in the
plus phase \cite[(1.7)--(1.10)]{ONGP}.

\begin{theorem}
$X_H$ is holomorphic on $|s|<1$, with $X_H(0)=1$, and $X_L$ is
holomorphic on $|s|>1$. The unit circle is a natural boundary for
each. The same boundary holds for $\chi_H$ and for every local branch
of $\chi_L$ continued from $s>1$ in the exterior. At every point of
their respective real temperature intervals the germs of $X_\nu$ and
$\chi_\nu$ are differentially transcendental over $\C(s)$.
\end{theorem}

Here differential transcendence means that there is no nonzero
polynomial relation among the function and finitely many derivatives
with coefficients rational in $s$. We use the dispersion variable
$E=s+s^{-1}$. Both temperature domains map to the slit plane
$\C\setminus[-2,2]$, with $s=0$ corresponding to infinity. In the
series normalization of Guttmann and Maillard \cite[(4)--(5)]{GM},
the high- and low-temperature series $\widetilde\chi_H(w)$ and
$\widetilde\chi_L(w)$ are also differentially transcendental over
$\C(w)$, where $w=1/(2E)$.

The proof turns on two estimates in the $E$-plane. At a boundary
energy $x$ we first bound the $j$th derivative of the whole sum over
particle numbers $N>2j+4$. This leaves finitely many integrals. At the
least resonant order $n$, the $n$-particle integral has a divergent
derivative of order $k$. All its collisions contribute with the same
phase. With $j=k$, the tail estimate transfers this divergence to
the full susceptibility. For differential transcendence we also need a
boundary interval on which any prescribed finite number of derivatives
have continuous limits.

\section{Site integrals and their sum}

For $q\in\T=\R/(2\pi\Z)$ let $z=e^{\ii q}$ and let $u=e^{-\gamma}$
be the root in the open unit disk of $u+u^{-1}=2(E-\cos q)$.
Write
\[
 K_{ij}=\frac{\sin((q_i-q_j)/2)}{\sinh((\gamma_i+\gamma_j)/2)},
 \qquad H(v,w)=\frac{v-w}{1-vw}.
\]
Subtracting the dispersion equations gives the single-valued square
\begin{equation}\label{pair}
 K_{ij}^2=-\frac{u_i u_j}{z_i z_j}
          \left(\frac{z_i-z_j}{1-u_i u_j}\right)^2
          = H(z_i,z_j)H(u_i,u_j).
\end{equation}
At $z_i z_j=1$ the middle expression gives the removable value of the product on the right.
For $N\ge1$, $A\in\Z_{\ge0}$ and $B\in\Z$, the real-temperature
spin-correlation expansion \cite[(3.13), (3.17)--(3.20), (4.1)--(4.2)]{ONGP}
gives the site coefficient
\begin{equation}\label{site}
 I_N(A,B;E)=\frac1{N!}\int_{\T^N}\prod_{i<j}K_{ij}^2
       \prod_i u_i^A z_i^B\,\frac{\dd q_i}{2\pi\sinh\gamma_i}.
\end{equation}
The ordered-state expansion also appears in \cite[Theorem~5.0]{PT}.
The connected correlation at $(a,b)$ is the sum of
$I_N(|a|,b;E)$ over positive odd $N$ in the high-temperature phase
and positive even $N$ in the low-temperature phase, multiplied by
$a_H=(1-s^4)^{1/4}/s$ and $a_L=(1-s^{-4})^{1/4}$, respectively.
The fourth roots tend to one at $s=0$ and $s=\infty$, respectively.

There are $N(N-1)/2$ pairs in \eqref{site}. Bounding them separately
would cost a constant raised to a power of order $N^2$. The following
bound for their product makes the particle series summable.

\begin{lemma}[Pair product on a common contour]\label{lem:pair-bound}
Suppose the angular contours and energies range over a compact family
on which the one-variable roots are holomorphic in a common neighborhood.
Assume that $|1-u_i u_j|\ge d>0$ for every pair of points on the same
contour, including coincident points. Then, for all $N$,
\[
 \left|\prod_{i<j}K_{ij}^2\right|\le (C N^{3/4})^N .
\]
The constant depends only on the family, the neighborhood and $d$.
\end{lemma}
\begin{proof}
For $D(v)=\prod_{i<j}H(v_i,v_j)^2$ the Cauchy identity gives
$D(v)=\det[1+H(v_i,v_j)]$. Indeed,
$1+H(v,w)=(1+v)(1-w)/(1-vw)$. On clearing the denominators of
$\det[(1-v_iw_j)^{-1}]$, alternation and degree give the two
Vandermonde factors, and induction at $v_Nw_N=1$ fixes the constant.

Translate all angles by the same complex number $t$. On a disk
$|t|<2r$ of fixed radius the roots stay holomorphic and retain a
reciprocal gap. Write $P(t)$ for the translated product of the squared
pairs. The middle expression of \eqref{pair} shows that $P$ is
holomorphic on this disk. Except at the apparent poles of the angular
determinant we have $P(t)^2=D(z(t))D(u(t))$.

In the root determinant every row has Euclidean norm at most $C\sqrt N$.
For an angular row $R_i(t)$ we estimate its mean over $|t|<r$,
with area normalized to one. The numerators of its entries are bounded.
The possible poles satisfy $z_i z_j e^{2\ii t}=1$ and are separated by
$\pi$ in $t$. Choose $r<\pi/8$. If there is a pole $t_0$ in $|t|<2r$,
factor $1-e^{2\ii(t-t_0)}$ there. On the smaller disk the entry is
bounded by $C(1+|t-t_0|^{-1})$, whose area mean is uniformly bounded
because $\int_{|w|<3r}|w|^{-1}\dd^2w=6\pi r$. In the absence of such
a pole the denominator on the smaller disk has a uniform gap by
compactness. Summing these entry bounds gives
$\langle\|R_i\|_2\rangle\le CN$.

Apply the submean inequality to $\log|P|$, Hadamard's inequality to
the two determinants and Jensen's inequality to each row norm. Then
\[
  2\log|P(0)|
   \le N\log(C\sqrt N)+\sum_{i=1}^N\log\langle\|R_i\|_2\rangle
   \le N\log(C\sqrt N)+N\log(CN).
\]
Taking square roots gives the asserted exponent. The factorial in
\eqref{site} makes $(CN^{3/4})^N/N!$ summable for every fixed $C$.
\end{proof}

The decay in both lattice coordinates follows from their interchange.
For real $E>2$, on circles $|v_i|=|w_i|=\rho<1$ close to one, integrate
\[
 \frac1{N!}\prod_{i<j}H(v_i,v_j)H(w_i,w_j)
 \prod_i\frac{v_i^A w_i^B}
 {E-(v_i+v_i^{-1}+w_i+w_i^{-1})/2}
 \frac{\dd v_i\dd w_i}{(2\pi\ii)^2v_iw_i},\qquad A,B\in\Z_{\ge0}.
\]
This is the double-circle representation of \cite[Appendix~1]{TW}
in even parity. For either parity, each $v_i$ integral has only the
interior root residue $u_i^A/\sinh\gamma_i$. Formula \eqref{pair}
then allows the $w_i$ circles to expand to unit radius, giving
\eqref{site}. Exchanging the circles, reflecting all $q_i$, and
continuing in $E$ gives $I_N(A,B)=I_N(B,A)=I_N(B,-A)$ for $A,B\ge0$.
On each compact subset of the slit plane the real-angle $|u|$ is
bounded by some $\rho_0<1$ and the other factors are bounded. Hence
$|I_N(|a|,b)|\le (C N^{3/4})^N\rho_0^{N\max(|a|,|b|)}/N!$, summably
in both sites and orders.

For $\sigma=\pm1$ put $H_N^\sigma=\{q:e^{\ii\sum q_i}=\sigma\}$,
$Z=\prod_i u_i$ and $\dd\mu=\prod_{i<N}\dd q_i/(2\pi)$.
Summing \eqref{site} with weight $\sigma^{a+b}$ gives
\begin{equation}\label{bulk}
 C_N^\sigma(E)=\sum_{a,b\in\Z}\sigma^{a+b}I_N(|a|,b;E)
 =\frac1{N!}\int_{H_N^\sigma}\frac{1+\sigma Z}{1-\sigma Z}
          \prod_{i<j}K_{ij}^2\prod_i\frac1{\sinh\gamma_i}\dd\mu .
\end{equation}
Summing $(\sigma Z)^{|a|}$ in $a$ gives the displayed quotient. For
the $b$ sum, take the total angle $\theta=\sum_i q_i$ as the last
integration variable. The remaining density has an absolutely
convergent Fourier series by the site bound. Its value at $\theta=0$
or $\pi$ according to the sign of $\sigma$ gives precisely the integral
over $H_N^\sigma$.
Thus $X_\nu=a_\nu F_\nu$, where $F_\nu(E)$ sums $C_N^+(E)$ in
the parity of the phase. Near $s=0$ the factors $u$, $1/\sinh\gamma$ and
$K_{ij}$ are uniformly $O(|s|)$, so $|C_N^+|\le3(C|s|)^{N^2}/N!$.
Since $C_1^+(E(s))=1/(E-2)=s/(1-s)^2$, the simple zero of $F_H(E(s))$
cancels the pole of $a_H$, giving $X_H(0)=1$.
Shifting all angles by $\pi$ and negating the roots gives
\begin{equation}\label{reflect}
 C_N^+(-E)=(-1)^N C_N^{(-1)^N}(E).
\end{equation}
In particular reflection in odd parity changes both signs in \eqref{bulk}.

\section{Contour bounds near the energy cut}

We use two deformations of the same one-angle contour. For a site
coefficient the center of the deformation is fixed by the lattice
direction. For a fixed particle integral its center is chosen to
preserve the total angle. In both cases the required decay comes from
the strict concavity of the boundary dispersion relation.

Fix $0<x<2$. At $E=x+\ii y$, $y>0$, write $\gamma=\ii\phi$ on
the physical lift. On the oscillatory interval $|r|<b=\arccos(x-1)$
the boundary value $g(r)=\arccos(x-\cos r)$ of $\phi$ satisfies
\[
 g'=-\frac{\sin r}{\sin g},\qquad
 g''=-\frac{x(1-\cos r\cos g)}{\sin^3g}<0 .
\]
Let $f$ be odd and periodic, equal to $r$ on $[-b,b]$ and linear
from $b$ to zero on $[b,\pi]$. For $c$ in a compact subinterval of
$(-b,b)$ put $\lambda=g'(c)$ and $q=r+\ii v(f(r)-c)$,
$-\pi\le r\le\pi$. Uniformly for small $v,y\ge0$,
\begin{equation}\label{phase}
 \Im(\lambda q-\phi(E,q))\ge a\{y+v(r-c)^2\},\qquad a>0 .
\end{equation}
On a closed interval inside $(-b,b)$ containing the centers, the
linear term at $y=0$ is
$v(r-c)(g'(c)-g'(r))\ge a_1v(r-c)^2$, with imaginary remainder
$O(v^3|r-c|^3)$. Increasing $y$ adds at least $a_2y$, since
$\Re(1/\sin\phi)>0$ there.

Near either endpoint $\sin r(f(r)-c)>a_3>0$, so
$W=E-\cos q$ satisfies $\Im W\ge y+a_4v$. Put $G=-\Im\phi$.
For $y+v>0$ we have $G>0$, and $W=\cos\phi$ gives
\[
 \frac{\Im W}{G}=\sin(\Re\phi)\,\frac{\sinh G}{G}\longrightarrow 0
 \qquad(W\longrightarrow1).
\]
By shrinking the endpoint neighborhoods we can therefore achieve
$G\ge M(y+a_4v)$ for any fixed $M$. Since $|\lambda\Im q|\le C v$
and $|r-c|\le2\pi$, taking $M$ large gives \eqref{phase} there.
On the remaining exterior intervals $G$ is bounded away from zero
already at $y=v=0$.

On the undeformed boundary contour the unit roots have arguments in
$[-b,0]$, with $b<\pi$. Two can be reciprocal only if both equal $1$.
Near that branch value $|u|=e^{-G}<1$ throughout the deformation
when $v+y>0$. Away from it compactness bounds both the one-body
and the pair denominators away from zero for small $v,y$. Fixing
one such $v>0$ now gives a positive gap everywhere on the graph.
For $x$ in a compact subinterval and centers a fixed distance from
$\pm b$, compactness gives the same gap, slightly decreased, after
small complex changes of $E$ and a common translation of the angles.
This neighborhood is fixed before the particle number is chosen.

\begin{proposition}[Site derivatives and the high-particle tail]\label{prop:site}
Fix a compact interval $I\subset(0,2)$ and an integer $j\ge0$.
Use lattice symmetry to put $A=\max(|a|,|b|)$ and $B=\min(|a|,|b|)$.
There are $y_0>0$ and $C_j$ such that for $x\in I$, $0\le y\le y_0$
and every $N\ge1$,
\begin{equation}\label{sites-bound}
 |\partial_E^j I_N(A,B;x+\ii y)|\le
 \frac{C_j^{\,N}N^{3N/4}}{N!}(1+A)^{j-N/2}.
\end{equation}
For either sign in \eqref{bulk} and either parity, the sum over
$N>2j+4$ therefore has continuous derivatives through order $j$ up
to $I$ from above.
\end{proposition}
\begin{proof}
For $A\ge1$, choose $g'(c)=B/A$. As $B/A$ runs from zero to one,
$c$ runs from zero to $-\arccos(x/2)$, where $g'(c)=1$.
These centers stay a positive distance from $\pm b$ for $x\in I$.
Fix a small positive $v$ in \eqref{phase}. On the deformed contour
the factor $u^A z^B=\exp(\ii A[(B/A)q-\phi])$ is bounded in modulus by
$\exp[-a' A(r-c)^2]$. The contour Jacobian and each one-body factor
are bounded, and the gaps just proved allow us to use
Lemma~\ref{lem:pair-bound} on a common complex neighborhood.

Keep this contour fixed while $E$ moves in a disk of radius
$\delta/(1+A)$ about $x+\ii y$. The derivative $\partial_E\phi$
is bounded on that neighborhood. Thus each factor $u^A z^B$ changes
in modulus by a factor at most $\exp(CA\delta/(1+A))$ relative to
its value at the center. The deformed integral agrees with the
physical one at positive heights by contour deformation, and hence
gives its germ on the disk. Uniformly on that disk its modulus is
at most
\[
 \frac{(C N^{3/4})^N}{N!}
 \left(\int_{-\pi}^{\pi}e^{-a'A(r-c)^2}\,\dd r\right)^N
 \le \frac{(C' N^{3/4})^N}{N!}(1+A)^{-N/2}.
\]
Cauchy's formula costs $j!\delta^{-j}(1+A)^j$. At the origin take
$c=0$ and a disk of fixed radius. This proves \eqref{sites-bound}.
The same disks furnish continuous limits at $y=0$.

There are $8A$ sites on the shell of radius $A\ge1$. For $N>2j+4$,
$A(1+A)^{j-N/2}\le (1+A)^{-3/2}$. The bound is now summable
separately over shells and particle numbers, uniformly up to the cut.
The same majorant applies to each lower derivative after increasing
$C_j$, which gives the asserted continuous traces.
\end{proof}

\paragraph{The remaining particle integrals.}
The one-particle term is $C_1^\sigma=1/(E-2\sigma)$.
Fix $N\ge2$ and $\sigma$, and exclude the values $x=1+\cos\theta$ with
$e^{\ii N\theta}=\sigma$. On $H_N^\sigma$ take
$c=N^{-1}\sum_i f(r_i)$. The excluded values rule out the two constant
tuples at $\pm b$, so $c$ stays in a compact subinterval of $(-b,b)$.
On the graphs
$q_i=r_i+\ii v(f(r_i)-c)$ one has $\sum_i\Im q_i=0$, so the
angular constraint is preserved.
Along this homotopy \eqref{phase} gives
\[
 -\log |Z|\ge a\bigl(Ny+v\textstyle\sum_i(r_i-c)^2\bigr).
\]
For $y>0$ this keeps $1-\sigma Z$ nonzero throughout the homotopy.
Together with the one-body and pair gaps, Stokes' theorem on the
smooth pieces of the graph gives \eqref{bulk} on the final cycle.
The shared faces have opposite orientations.
At fixed small $v=\epsilon>0$, the only possible poles at $y=0$ are
the constant configurations satisfying
\begin{equation}\label{res}
 e^{\ii N q_0}=e^{\ii N g(q_0)}=\sigma,\qquad
 x=\cos q_0+\cos g(q_0),\qquad |q_0|<b .
\end{equation}
Choose disjoint neighborhoods of these tuples and a nonnegative cutoff
in each, equal to one on a smaller neighborhood. On the remaining
part of the fixed graph all denominators have a gap, so its integral
is holomorphic across $x$. In a collision chart write $r_i=q_0+v_i$. The
constraint gives $\sum_i v_i=0$ and $f(r_i)=r_i$, so the center
of the deformation is exactly $q_0$. Hence $q_i=q_0+L v_i$ with
$L=1+\ii\epsilon$ for every label.

Put $\eta=E-x$, $D=1-\sigma Z$, $Q=\sum_i v_i^2$ and
$\Delta(v)=\prod_{i<j}(v_i-v_j)$. Write $\psi=g(q_0)$.
Since $\phi_E=-1/\sin\psi$ at the center and $\sum_i v_i=0$,
Taylor expansion gives
\begin{equation}\label{expansion}
 D=-\ii(d\eta+hL^2 Q)+O(|\eta|^2+|\eta|\,|v|+|v|^3),\qquad
 d=\frac{N}{\sin\psi}>0,\qquad h=-\frac{g''(q_0)}2>0 .
\end{equation}
For $\Im\eta\ge0$, $|d\eta+hL^2Q|\ge c(|\eta|+Q)$ because
$0<\arg L^2<\pi$. This estimate and the gap off the charts identify
the fixed-graph integral with the physical germ on an upper half-disk
by continuation from $\eta=\ii y$. Each squared pair supplies
$(v_i-v_j)^2$ times an analytic unit. The density in
$\prod_{i<N}\dd v_i$ is $\Delta^2 B(\eta,v)/D(\eta,v)$, where
$B(0,0)\ne0$ since the one-body factors, the graph Jacobian and
$1+\sigma Z$ are nonzero there. In polar coordinates $r=\sqrt Q$
the volume element of the $(N-1)$-dimensional plane contributes
$r^{N-2}$ and the squared Vandermonde contributes $r^{N(N-1)}$.
Thus the radial measure is proportional to $r^{N^2-2}\dd r$ after
angular integration. The angular constant is positive since tuples
of distinct coordinates occupy a relatively open set of the unit
sphere. In particular, all poles $D^{-m}$ with
$N^2-2-2m>-1$ have a common integrable majorant for $\eta$ in a
sufficiently small upper half-disk at fixed $N,m$.

\section{The first resonant order}

Take $\sigma=\pm1$ in odd parity and $\sigma=1$ in even parity.
Choose $0<\alpha<\pi/2$ with $e^{\ii N\alpha}=\sigma$ for some $N$
of the chosen parity. The mesh of these angular grids is $2\pi/N$,
so the energies $x=2\cos\alpha$ are dense in $(0,2)$ even after
omitting $x=1$. Since $g(\alpha)=\alpha$, the constant tuple at
$\alpha$ satisfies \eqref{res}. Let $n$ be the least resonant
particle number of this parity at $x$.

\begin{proposition}[Noncancellation at the first resonance]\label{prop:resonance}
For such a sign, parity and energy $x\ne1$, let
$F^\sigma(E)=\sum_N C_N^\sigma(E)$ in the given parity and set
$k=\lceil(n^2-3)/2\rceil$. There is a ray $E=x+L^2t$ in the upper
half-plane along which
\[
 \bigl|\partial_E^k F^\sigma(x+L^2t)\bigr|\longrightarrow\infty,
 \qquad
 \partial_E^j F^\sigma(x+L^2t)=O(1)\quad(0\le j<k),
 \qquad t\downarrow0 .
\]
\end{proposition}
\begin{proof}
We first exclude the branch endpoints, at every particle number.
No such $x$ is a value $1+\cos\theta$ for a root-of-unity
angle $\theta$. Indeed $x=e^{\ii\alpha}+e^{-\ii\alpha}$ is an algebraic
integer. If $u=\cos\theta=x-1$ then $u$ is also an algebraic integer.
For $m\theta\in2\pi\Z$ its conjugates are roots of $T_m(X)-1$, where
$T_m(\cos t)=\cos(mt)$. All those roots lie in $[-1,1]$.
If any conjugate is $1$ or $-1$, irreducibility gives $u=1$ or
$u=-1$. Otherwise the product of the conjugates is an integer of
absolute value less than one, so $u=0$. These cases give $x=2,0,1$,
all excluded.

Choose a single contour displacement $\epsilon>0$ small enough for all
orders $N\le 2k+4$, and put $L=1+\ii\epsilon$.
At a collision of order $n$ the phase can be read off from the
original integral \eqref{bulk}. Put $\psi=g(q_0)\in(0,\pi)$ and
$p=n(n-1)/2$. Dividing a squared pair by $(v_i-v_j)^2$ gives the
center value $L^2/[4(\ii\sin\psi)^2]$. The $n$ one-body factors
give $(\ii\sin\psi)^{-n}$ and the constrained angular Jacobian is
$L^{n-1}$. With the positive Haar normalization and
$1+\sigma Z=2$, the center value in the density $\Delta^2 B/D$ is
\[
 B(0,0)=
 \frac{2}{n!(2\pi)^{n-1}}\,
 \frac{L^{\,2p+n-1}}{4^p(\ii\sin\psi)^{2p+n}}
 = b_{q_0}\,\ii^{-n^2}L^{\,n^2-1},\qquad b_{q_0}>0 .
\]
Thus the center values differ only by the positive factors $b_{q_0}$.

Take $\eta=L^2t$, $t\downarrow0$. Equation~\eqref{expansion} gives
$D=-\ii L^2(dt+hQ)(1+O(t+|v|))$ on shrinking the chart, since
$dt+hQ$ is comparable to $t+Q$. The term of pole order $k+1$ in
the $k$th derivative of $B/D$ is
\[
 (-1)^k k!\,B\,(\partial_\eta D)^k/D^{k+1}.
\]
Every other term has pole order at most $k$, and its integral is
bounded since $n^2-2-2k>-1$.
In the displayed term $\partial_\eta D=-\ii d+O(t+|v|)$. Apart from
a relative error $O(t+|v|)$ the term is therefore
\[
 \underbrace{\frac{(-1)^k k!\,\ii^{-n^2}L^{n^2-1}(-\ii)^k}
                    {(-\ii L^2)^{k+1}}}_{\displaystyle\omega}\,
 \frac{b_{q_0}d^k}{(dt+hQ)^{k+1}} .
\]
The complex number $\omega\ne0$ is the same in every chart. Multiply
by $\overline\omega/|\omega|$ and shrink the finitely many charts
so the relative error is at most $1/2$. The real part is then bounded
below by $c(t+Q)^{-k-1}$. Integrating against $\Delta^2$ on
$\sqrt t<r<\delta$ gives a positive constant times
\[
 \int_{\sqrt t}^\delta r^{\,n^2-2-2(k+1)}\dd r .
\]
The exponent is $-1$ for odd $n$ and $-2$ for even $n$. In either
case the rotated real parts add and diverge. Derivatives of order
$j<k$ have pole order at most $k$, so they remain bounded.

It remains to put back the other particle orders. Every $N<n$ in the
chosen parity is nonresonant and continues across $x$. For
$n<N\le2k+4$ the endpoint test still applies. At a regular collision
the worst pole in the $k$th derivative is now integrable, because
$N^2-2-2(k+1)>-1$. Finally, Proposition~\ref{prop:site} bounds the
derivatives through order $k$ of the infinite tail $N>2k+4$ on a
closed upper half-disk about $x$.
These three contributions are bounded on the same ray.
\end{proof}

\section{Natural boundary and differential transcendence}

\paragraph{The natural boundary.}
For positive $x$ use Proposition~\ref{prop:resonance} with $\sigma=1$.
For negative energies use
\eqref{reflect} with $\sigma=(-1)^n$. Complex conjugation supplies
the other side of the slit. The map $E=s+s^{-1}$ is locally
invertible at the resulting dense set on $|s|=1$. The factors
$a_\nu$ and every continued branch of $\arsinh(s)/(2J)$ are analytic
units there. By the chain and product rules the $k$th derivative of
each of $F_\nu$, $X_\nu$, and $\chi_\nu$ diverges along the corresponding
approach from within its domain, with all lower derivatives bounded. Density proves the
natural-boundary assertions, including at the omitted axes.

\paragraph{The boundary trace of a differential relation.}
Suppose one of these functions, in the coordinate $z=E$ for $F_\nu$
or $z=s$ for $X_\nu$ or $\chi_\nu$, satisfies a differential
relation. After clearing denominators take $P(z,Y,\ldots,Y^{(r)})=0$
of least order and then least degree in its highest derivative.
Its separant $S=P_{Y_r}(z,Y,\ldots,Y^{(r)})$ is nonzero as a germ,
since the formal derivative is a nonzero polynomial of smaller degree
in $Y_r$ and so cannot give a relation. Continue the relation from the real interval to an
upper energy collar in $0<\Re E<2$, fixing the local inverse $s(E)$
and the physical branches. Delete
from $(0,2)$ the finitely many endpoint and two-cosine grid values
of orders $N\le2r+4$. On a remaining open subinterval the finitely
many integrals of those orders continue across the cut. On each compact
subinterval Proposition~\ref{prop:site} gives continuous one-sided
limits for the derivatives through $r$ of their complement. After
composition with the local chart, $S$ thus has a continuous trace there.

Each bounded range of particle numbers supplies only finitely many
two-cosine grid values in \eqref{res}. Thus the selected $x$ with
$k>r$ are dense in this subinterval. At each such point, differentiate the
relation $k-r$ times. The chain rule gives
\[
 S\,Y^{(k)}+R(z,Y,\ldots,Y^{(k-1)})=0 .
\]
The remainder stays bounded along the chosen ray, whereas
$|Y^{(k)}|\to\infty$. The trace of $S$ is therefore zero at each
of these points and, by continuity, on the whole subinterval.
Schwarz reflection in $E$ across this zero trace and
the identity theorem force $S$ to vanish identically, a contradiction.
Finally, $X_\nu=a_\nu F_\nu$ gives
$\widetilde\chi_\nu(w)=F_\nu(1/(2w))$ in the normalization above.
A differential relation for either germ at $w=0$ would yield one
for $F_\nu(E)$ by the chain rule, since $w=1/(2E)$ is birational.

\begingroup\small
\begin{thebibliography}{9}
\setlength{\itemsep}{3pt}\setlength{\parsep}{0pt}
\bibitem{GE} A. J. Guttmann and I. G. Enting,
\emph{Solvability of some statistical mechanical systems},
Phys. Rev. Lett. \textbf{76}(3) (1996), 344--347.
\href{https://doi.org/10.1103/PhysRevLett.76.344}{doi:10.1103/PhysRevLett.76.344}.
\bibitem{Nickel} B. Nickel, \emph{On the singularity structure of the
2D Ising model susceptibility}, J. Phys. A \textbf{32}(21) (1999),
3889--3906. \href{https://doi.org/10.1088/0305-4470/32/21/303}{doi:10.1088/0305-4470/32/21/303}.
\bibitem{GJMP} A. J. Guttmann, I. Jensen, J.-M. Maillard and J. Pantone,
\emph{Is the full susceptibility of the square-lattice Ising model a
differentially algebraic function?}, J. Phys. A \textbf{49}(50) (2016),
504002. \href{https://doi.org/10.1088/1751-8113/49/50/504002}{doi:10.1088/1751-8113/49/50/504002}.
\bibitem{GM} A. J. Guttmann and J.-M. Maillard,
\emph{Automata and the susceptibility of the square lattice Ising model
modulo powers of primes}, J. Phys. A \textbf{48}(47) (2015), 474001.
\href{https://doi.org/10.1088/1751-8113/48/47/474001}{doi:10.1088/1751-8113/48/47/474001}.
\bibitem{ONGP} W. P. Orrick, B. Nickel, A. J. Guttmann and J. H. H. Perk,
\emph{The susceptibility of the square lattice Ising model: New developments},
J. Stat. Phys. \textbf{102}(3--4) (2001), 795--841.
\href{https://doi.org/10.1023/A:1004850919647}{doi:10.1023/A:1004850919647}.
\bibitem{PT} J. Palmer and C. Tracy, \emph{Two-dimensional Ising
correlations: Convergence of the scaling limit}, Adv. Appl. Math.
\textbf{2}(3) (1981), 329--388.
\href{https://doi.org/10.1016/0196-8858(81)90010-5}{doi:10.1016/0196-8858(81)90010-5}.
\bibitem{TW} C. A. Tracy and H. Widom, \emph{On the singularities in the
susceptibility expansion for the two-dimensional Ising model},
J. Stat. Phys. \textbf{156}(6) (2014), 1125--1135.
\href{https://doi.org/10.1007/s10955-014-1061-4}{doi:10.1007/s10955-014-1061-4}.
\end{thebibliography}
\endgroup
\end{document}
